problem:  
Let \( (M,g) \) be a compact oriented Riemannian surface, and \( (N,h) \) a compact Riemannian manifold.  
For each \( n \), let \( G_n = (V_n,E_n,F_n) \) be a uniformly shape-regular geodesic triangulation of \( M \) with mesh size \( h_n \to 0 \).  

For each \( n \), let \( f_n : V_n \to N \) be a **discrete generalized harmonic map**, i.e. a critical point of  
\[
E_n(f_n) = \frac{1}{2}\sum_{(i,j)\in E_n} w_{ij}^{(n)} \Phi\!\left(d_N^2(f_n(i), f_n(j))\right),
\]
where \( \Phi:[0,\infty)\to[0,\infty) \) is smooth, strictly convex with \( \Phi'(0)=1 \), and the intrinsic weights \(w_{ij}^{(n)}\) are chosen so that \(E_n\) $\Gamma$‑converges to  
\[
E(f) = \tfrac{1}{2}\int_M |df|_g^2\,d\mathrm{vol}_g .
\]
Assume \( f_n \) converges weakly in \(H^1\) to a smooth harmonic map \(f:M\to N\).

For each face \(T=(v_i,v_j,v_k)\in F_n\), consider the piecewise-geodesic loop in \(N\) joining 
\( f_n(v_i)\to f_n(v_j)\to f_n(v_k)\to f_n(v_i) \) sequentially along minimizing geodesics.  
Let \( P_T^{(n)}: T_{f_n(v_i)}N \to T_{f_n(v_i)}N\) be the **holonomy operator** obtained by parallel transport along this loop.  
Define the **discrete curvature probe** by
\[
\mathcal{K}_T^{(n)} := \frac{1}{A_T^{(n)}} \log(P_T^{(n)}) \;\in\; \mathfrak{so}(T_{f_n(v_i)}N),
\]
where \(A_T^{(n)}\) is the area of the domain triangle in \(M\),  
\(\log\) denotes the principal Riemannian logarithm (well-defined for small loops),  
and the norm \(\|\cdot\|\) is taken with respect to \(h\).  
For smooth \(f\), infinitesimal holonomy expansions yield
\(\log(P_T^{(n)}) \approx A_T^{(n)}\,R_{f_n(v_i)}(df_n(e_1),df_n(e_2))\),  
so that \( \mathcal{K}_T^{(n)} \) approximates the curvature 2‑form of \(N\) along the image plane of \(df\).

Define the measure on \(M\)
\[
\mu_n := \sum_{T\in F_n} A_T^{(n)} \big\| \mathcal{K}_T^{(n)} \big\|\, \delta_{p_T},
\]
with \(p_T\) the barycenter of \(T\).  (The norm of \(\log(P_T^{(n)})\) is invariant under change of the base vertex of the loop up to \(O(h_n^3)\), ensuring consistency.)

**Open problem:**  
Does the sequence \(\mu_n\) converge (perhaps along a subsequence) in the weak‑* sense of measures to  
\[
\mu = c_*\, |K_N(df_p(e_1)\wedge df_p(e_2))|\, |\!df_p(e_1)\wedge df_p(e_2)\!|\,d\mathrm{vol}_g(p),
\]
for some universal nonzero constant \(c_*\) depending only on the local geometric normalization of the triangulation, not on \(M,N,f\) or \(\Phi\)?  

Equivalently: can the curvature measure of the target manifold \(N\) emerge as the continuum limit of rescaled holonomy‑defect measures induced by \emph{any} sequence of discrete generalized harmonic maps, without explicitly encoding curvature in the discrete variational energy?

---

Why is it a "good" problem:  
This problem isolates a geometrically plausible yet analytically uncharted phenomenon—**curvature emergence through the asymptotic holonomy structure of discrete variational maps**.  

It achieves rigorous geometric coherence:
- Curvature genuinely enters through holonomy (a second‑order effect).
- The base‑point ambiguity and scaling are now clarified.
- The dependence on the energy function \(\Phi\) is removed, leaving an intrinsic geometric constant \(c_*\).

Conceptually, the problem bridges discrete harmonic analysis, Γ‑convergence theory, and Riemannian curvature geometry via parallel transport.  
A positive answer would reveal that curvature information of \(N\) is already implicitly coded in the discrete holonomy fields generated by variationally defined harmonic maps—a deep unification of first‑order discrete analysis and second‑order geometric structure.  
A negative answer would delineate the fundamental barrier delineating what discrete variational principles can or cannot reconstruct about continuous curvature.

The formulation is compact yet substantive, requiring mastery of discrete connection theory, measure convergence, and curvature expansions—precisely the interplay of simplicity, depth, and cross‑field relevance that defines a “good” research‑level mathematical problem.